\subsection{Evidence extraction and synthesis}
Use one study record per source based on the shared template. One author
extracts and the other verifies against the source. Record citation key
and page/table/figure locations, task, domain, datasets, generator, splits,
training regime, requirements (RQ1), metrics and limitations (RQ2), and
assessment prerequisites (RQ3). Record feature networks and checkpoints,
preprocessing, sample sizes, baselines, subgroup definitions, uncertainty,
code availability and limitations where reported.

Compare evidence within compatible tasks and settings. Separate source
claims, reported results and our interpretation. Do not pool scores from
incompatible feature extractors, datasets, preprocessing or training
budgets. Use narrative synthesis and a requirement--metric--resource
mapping. Quantitative pooling requires an explicit comparability rationale.
Report missing evidence, conflicting results and strength of support.

\subsection{Review protocol versus experimental assessment}
This paper currently plans a literature review, not a new training
experiment. Section~\ref{sec:rq3} synthesises assessment prerequisites.
If an experiment is added, specify its protocol before running it:
independent real train/validation/test splits; real-only, synthetic-only
and mixed-data baselines; controlled architectures and training budgets;
multiple seeds; task and subgroup metrics; uncertainty estimates; and
safeguards against test-set tuning.
